\documentclass[a4paper,12pt]{article}
\usepackage{cmap}
\usepackage[T2A]{fontenc}
\usepackage[utf8]{inputenc}
\usepackage[russian]{babel}
\usepackage{amsmath,amssymb,mathtools}
\usepackage{helvet}
% Helvetica has no T2A glyphs; use its Cyrillic sans-serif counterpart.
\DeclareFontFamily{T2A}{phv}{}
\DeclareFontShape{T2A}{phv}{m}{n}{<->ssub*cmss/m/n}{}
\DeclareFontShape{T2A}{phv}{b}{n}{<->ssub*cmss/bx/n}{}
\DeclareFontShape{T2A}{phv}{m}{it}{<->ssub*cmss/m/sl}{}
\renewcommand{\familydefault}{\sfdefault}
\usepackage{icomma}
\usepackage[expansion=false]{microtype}
\usepackage[a4paper,margin=19mm]{geometry}
\usepackage{tikz}
\usetikzlibrary{calc,arrows.meta,positioning,shapes.geometric}
\usepackage{pgfplots}
\pgfplotsset{compat=1.18}
\usepackage{longtable,booktabs,enumitem,parskip,fancyhdr,setspace,xcolor,hyperref}
\definecolor{accent}{HTML}{2F6F6D}
\definecolor{darkaccent}{HTML}{1F4E4D}
\definecolor{textgray}{HTML}{333333}
\hypersetup{pdftitle={Модуль и кусочные графики — Марина Селиверстова},pdfauthor={Лёвин Артём Александрович},colorlinks=true,linkcolor=darkaccent}
\onehalfspacing
\setlength{\parindent}{0pt}
\setlength{\headheight}{15pt}
\setlist[itemize]{leftmargin=5mm,itemsep=3pt,topsep=3pt}
\setlist[enumerate]{leftmargin=7mm,itemsep=9pt,topsep=4pt}
\renewcommand{\arraystretch}{1.3}
\pagestyle{fancy}\fancyhf{}
\fancyhead[L]{\small Модуль и кусочные графики}
\fancyhead[R]{\small 05.10.2026}
\fancyfoot[C]{\small\thepage}
\newcommand{\topic}[1]{\section*{\color{darkaccent}#1}}
\newcommand{\key}[1]{\textbf{\color{darkaccent}#1}}
\tikzset{startstop/.style={ellipse,draw=accent,align=center,minimum width=40mm,minimum height=11mm},action/.style={rectangle,rounded corners=2mm,draw=accent,align=center,text width=48mm,minimum height=17mm},decision/.style={diamond,aspect=2.3,draw=darkaccent,align=center,inner sep=2mm,text width=30mm},arr/.style={-{Latex[length=3mm]},draw=textgray,line width=.8pt}}
% Lesson date: 2026-10-05. Preparation date: 2026-10-05 (Europe/Moscow).
\begin{document}
\thispagestyle{empty}
\begin{center}
\vspace*{22mm}
{\Large\color{textgray}Математика · ОГЭ}\par
\vspace{11mm}
{\Huge\bfseries\color{darkaccent}Чек-лист}\par
\vspace{8mm}
{\LARGE\bfseries Модуль и кусочные графики}\par
\vspace{10mm}
{\large Подмодульное выражение, скобки, области ветвей}\par
\vspace{12mm}
{\large 5 октября 2026 года}\par
\vfill
{\large Лёвин Артём Александрович}\par
{\large эксклюзивно для Марины Селиверстовой}\par
\vspace{15mm}
\begin{tikzpicture}
\draw[accent!65,line width=1pt] (0,0) .. controls (3,.5) and (7,-.5) .. (11,0);
\end{tikzpicture}
\vspace{8mm}
\end{center}
\clearpage
\topic{1. Главная идея: знак всего выражения}
\key{Модуль числа} — расстояние от этого числа до нуля. Его значение всегда неотрицательно. Для любого действительного выражения $u$:
\[
|u|=\begin{cases}u,&u\ge0,\\-u,&u<0.\end{cases}
\]
В записи $|x-6|$ подмодульное выражение — $x-6$. Условия двух случаев: $x-6\ge0$ и $x-6<0$. Для линейного выражения их удобно записать как $x\ge6$ и $x<6$.

\key{Чек-лист раскрытия модуля}
\begin{itemize}
\item Выделите всё выражение между знаками модуля.
\item В условии случая сохраните это выражение без изменения знаков.
\item При неотрицательном выражении уберите модуль.
\item При отрицательном выражении замените модуль минусом и скобками вокруг всего выражения.
\item Сохраните множители, слагаемые и знаки снаружи модуля.
\end{itemize}
Нуль включаем в первую ветвь: две области вместе охватывают все допустимые значения и не пересекаются.

\topic{2. «Галка» и её сдвиг}
\[
y=|x-3|=\begin{cases}x-3,&x\ge3,\\-x+3,&x<3.\end{cases}
\]
Вершина имеет координаты $(3;0)$: график сдвинут вправо на $3$. У функции $y=|x+3|$ вершина $(-3;0)$: сдвиг влево на $3$. Быстрая проверка: подмодульное выражение в вершине равно нулю.
\begin{center}
\begin{tikzpicture}
\begin{axis}[width=12cm,height=7.5cm,axis lines=middle,xmin=-2,xmax=8,ymin=-1,ymax=6,unit vector ratio=1 1,samples=201,clip=false,xtick={0,3,6},ytick={0,3},xlabel={$x$},ylabel={$y$}]
\addplot[accent,very thick,domain=-2:3]{3-x};
\addplot[accent,very thick,domain=3:8]{x-3};
\addplot[only marks,mark=*,accent] coordinates {(3,0)};
\node[anchor=north,fill=white,inner sep=2pt] at (axis cs:3,-.15) {$(3;0)$};
\end{axis}
\end{tikzpicture}
\end{center}
\clearpage
\topic{3. Внешний минус и скобки}
Для функции $y=x^2-|x|-6$ меняется только часть $|x|$:
\[
y=\begin{cases}x^2-x-6,&x\ge0,\\x^2+x-6,&x<0.\end{cases}
\]
Во второй ветви: $x^2-(-x)-6=x^2+x-6$.

Теперь модуль охватывает выражение $x-6$:
\[
y=x^2-|x-6|.
\]
\begin{align*}
x-6\ge0:&\quad y=x^2-(x-6)=x^2-x+6;\\
x-6<0:&\quad y=x^2-\bigl(-(x-6)\bigr)=x^2+x-6.
\end{align*}
Поэтому
\[
y=\begin{cases}x^2-x+6,&x\ge6,\\x^2+x-6,&x<6.\end{cases}
\]
\key{Самопроверка знаков.} При $x=0$ исходная функция равна $-6$, и вторая ветвь также даёт $-6$. При $x=6$ обе формулы дают $36$.

\topic{4. Квадратный трёхчлен под модулем}
\[
y=4-|x^2-3x|=\begin{cases}4-x^2+3x,&x^2-3x\ge0,\\4+x^2-3x,&x^2-3x<0.\end{cases}
\]
Здесь достаточно записать условия относительно \emph{всего} трёхчлена. Для раскрытия модуля решать эти неравенства относительно $x$ пока не требуется.

Если модуль охватывает всю формулу, получаем:
\[
y=|x^2-x-6|=\begin{cases}x^2-x-6,&x^2-x-6\ge0,\\-x^2+x+6,&x^2-x-6<0.\end{cases}
\]
Минус перед скобками меняет знак каждого слагаемого. Обе функции определены при всех действительных $x$.
\clearpage
\topic{5. Границы модуля определяют ветви}
Сравните две функции:
\[
f(x)=|x|(x-6)+2,\qquad g(x)=x|x-6|+2.
\]
У первой функции под модулем только $x$:
\[
f(x)=\begin{cases}x(x-6)+2,&x\ge0,\\-x(x-6)+2,&x<0.\end{cases}
\]
У второй функции под модулем $x-6$:
\[
g(x)=\begin{cases}x(x-6)+2,&x\ge6,\\-x(x-6)+2,&x<6.\end{cases}
\]
Внешнее слагаемое $+2$ сохраняется в обеих ветвях. При $x=3$ значения различаются: $f(3)=-7$, $g(3)=11$. Одинаковые формулы ветвей могут применяться на разных областях.

\topic{6. Как собрать кусочный график}
\begin{enumerate}[itemsep=3pt]
\item Запишите формулу и область каждой ветви.
\item Постройте каждую ветвь только на её области.
\item Проверьте граничное значение: точка принадлежит ветви при нестрогом неравенстве.
\item Для формул, полученных раскрытием модуля, проверьте совпадение значений там, где подмодульное выражение равно нулю.
\end{enumerate}
Например, у $y=|x-3|$ обе формулы дают $0$ при $x=3$; вершина входит в график. В произвольной кусочной функции значения соседних формул на границе могут различаться.

\key{Нули функции} — значения аргумента, при которых $y=0$. Решайте уравнение для каждой ветви и проверяйте принадлежность найденного корня её области.

\key{Частые ошибки}
\begin{itemize}
\item Подменять условие $x-6\ge0$ условием $x\ge0$.
\item Терять внешний минус или раскрывать его без скобок.
\item Менять знак самого выражения в условии случая.
\item Продолжать ветвь за пределы её области или забывать граничную точку.
\end{itemize}
\clearpage
\topic{Алгоритм: раскрытие одного модуля}
\begin{center}
\begin{tikzpicture}[node distance=14mm,font=\small]
\node[startstop] (start) {Начало};
\node[action,below=of start] (read) {Выделите всё выражение $u$ под модулем};
\node[decision,below=of read] (cond) {$u\ge0$?};
\node[action,below left=18mm and 13mm of cond] (yes) {Ветвь $u\ge0$:\\замените $|u|$ на $u$};
\node[action,below right=18mm and 13mm of cond] (no) {Ветвь $u<0$:\\замените $|u|$ на $-(u)$};
\node[action,below=54mm of cond,text width=70mm] (outer) {В каждой ветви сохраните внешние знаки и множители; раскройте скобки};
\node[action,below=of outer,text width=70mm] (check) {Запишите обе формулы со своими условиями; проверьте значение при $u=0$};
\node[startstop,below=of check] (end) {Кусочная запись готова};
\draw[arr] (start)--(read);
\draw[arr] (read)--(cond);
\draw[arr] (cond.west)-|node[pos=.25,above]{да}(yes.north);
\draw[arr] (cond.east)-|node[pos=.25,above]{нет}(no.north);
\draw[arr] (yes.south)|-(outer.west);
\draw[arr] (no.south)|-(outer.east);
\draw[arr] (outer)--(check);
\draw[arr] (check)--(end);
\end{tikzpicture}
\end{center}
\clearpage
\topic{Тренировка · Путь А · 10 заданий}
Запишите развёрнутые преобразования. В заданиях 3--9 представьте функцию без знака модуля с условиями для каждой ветви. Квадратные неравенства в заданиях 6--7 решать не требуется.
\begin{enumerate}
\item Вычислите $|{-5}|$, $|0|$ и $|3-7|$.
\item Найдите вершину графика $y=|x+2|$. Укажите его сдвиг относительно графика $y=|x|$.
\item $y=|x-4|$.
\item $y=x^2-|x|-4$.
\item $y=x^2-|x-5|$.
\item $y=3-|x^2-2x|$.
\item $y=|x^2+x-2|$.
\item $y=|x|(x-4)+1$.
\item $y=x|x-4|+1$.
\item Для $y=|x+1|-2$ запишите кусочную формулу, найдите вершину и нули. Постройте график и проверьте значение в точке смены ветвей.
\end{enumerate}
\vfill
\key{Перед сверкой:} у каждой формулы есть условие; нуль включён ровно в одну ветвь; внешние знаки сохранены; корни прошли проверку по области ветви.
\clearpage
\topic{Ответы · задания 1--10}
\begin{longtable}{@{}r p{.92\linewidth}@{}}
\toprule
№ & Ответ \\
\midrule
\endhead
1 & $5;\ 0;\ 4$.\\[3mm]
2 & $(-2;0)$; сдвиг влево на $2$.\\[3mm]
3 & $\displaystyle y=\begin{cases}x-4,&x\ge4,\\-x+4,&x<4.\end{cases}$\\[7mm]
4 & $\displaystyle y=\begin{cases}x^2-x-4,&x\ge0,\\x^2+x-4,&x<0.\end{cases}$\\[7mm]
5 & $\displaystyle y=\begin{cases}x^2-x+5,&x\ge5,\\x^2+x-5,&x<5.\end{cases}$\\[8mm]
6 & $\displaystyle y=\begin{cases}3-x^2+2x,&x^2-2x\ge0,\\3+x^2-2x,&x^2-2x<0.\end{cases}$\\[7mm]
7 & $\displaystyle y=\begin{cases}x^2+x-2,&x^2+x-2\ge0,\\-x^2-x+2,&x^2+x-2<0.\end{cases}$\\[7mm]
8 & $\displaystyle y=\begin{cases}x^2-4x+1,&x\ge0,\\-x^2+4x+1,&x<0.\end{cases}$\\[7mm]
9 & $\displaystyle y=\begin{cases}x^2-4x+1,&x\ge4,\\-x^2+4x+1,&x<4.\end{cases}$\\[7mm]
10 & $\displaystyle y=\begin{cases}x-1,&x\ge-1,\\-x-3,&x<-1.\end{cases}$
\par Вершина $(-1;-2)$; нули $-3$ и $1$. График — две полупрямые с общей включённой вершиной; обе формулы на границе дают $-2$.\\[3mm]
\bottomrule
\end{longtable}
\end{document}
